\exnumber{00}
\extitle{UX VII - Ex 1 - Stratégie UX}
\exfiles{ex00/fichier Figma et PDF, nomme Name\_FirstName\_DeliverableName\_Date\_VersionNumber}
\exauthorize{Ceux nommés dans l'exercice}
\makeheaderfiles

\textbf{Contexte :}

Elsa est couturière. Elle fabrique des maillots de bain sur mesure via son compte Instagram. Ça marche bien, sa communauté grandit ainsi que le volume de demandes de clients.

Cependant, elle perd beaucoup de temps à chaque demande. Pour finaliser une commande, elle doit demander la forme du maillot, la taille, les préférences de couleur et le délai de livraison souhaité. En plus de cela, elle doit s'assurer d'avoir assez de temps pour coudre une pièce  ---  chaque pièce nécessite 6 heures.

Pour l'instant, elle ne travaille qu'avec des messages privés mais elle aimerait passer moins de temps sur cette partie.

Vous lui proposez de l'aider à concevoir une solution pour mieux gérer son temps et son atelier. Vous avez 4 jours de travail.

\textbf{Instructions :}

Maintenant que vous avez découvert et pratiqué de nombreux outils du processus de design, organisez un planning pour les 4 prochains jours de travail.

Sélectionnez au moins un outil pour chaque phase, nommez-le, et tenez compte du temps que chaque phase peut nécessiter.

\begin{itemize}
\item Empathie / Recherche utilisateur
\item Définir
\item Problem Statement
\item Idéation
\item Prototype
\item Test
\end{itemize}

\textbf{Livrables :}

\begin{itemize}
\item Un calendrier, en PDF ou sur Figma, du projet, comprenant les jours ouvrés à venir, chaque phase du processus, les principaux outils que vous prévoyez d'utiliser pour chaque phase et le temps que vous prévoyez de consacrer à chaque outil.
\item Déposez le PDF ou le document Figma sur Github, nommé comme suit : « Nom\_Prénom\_NomDuLivrable\_Date\_NumeroDeVersion ». Par exemple, cela peut être nommé « Doe\_John\_MidFiPrototype\_05242024\_V1 ».
\item Ajoutez un titre dans le document.
\item Soignez l'aspect global du document. Il doit être clair, simple et lisible. Vous pouvez regarder canva.com, sans surcharger le dessin de détails.
\end{itemize}

\textbf{Ressource :}

\begin{itemize}
\item Un récapitulatif de tous les outils qui ont été pratiqués \url{https://www.figma.com/file/x0xfoz9TaEDiAfGPFqKJcs/RECAP-UX%2FUI-Program?node-id=0%3A1}. Concentrez-vous sur la partie UX !
\end{itemize}

\textbf{Citations :}

\begin{itemize}
\item Le design a besoin d'échouer. L'échec est même une étape nécessaire, mais idéalement il doit survenir avant le lancement d'un produit, pendant les phases de prototype et de test
\item « La UX Strategy se situe à l'intersection de l'UX et du business. Elle offre bien plus de chances de créer des produits à succès. Elle permet aux équipes de voir la « big picture » pour atteindre les objectifs business dans des conditions d'incertitude » Jamie Levy | UX Strategy: How to Devise Innovative Digital Products That People Want
\item La UX Strategy est la méthode par laquelle vous validez que votre solution résout un problème pour de vrais clients dans un marché dynamique, car le marché évolue constamment.
\end{itemize}

